\begin{pastexam}[label={px:ntu-108-ee-dsb-05}]{臺大電機所 108 學年度（B）第 5 題}
（是非；答對 4 分，答錯倒扣 4 分）

Computers handle postfix notation more easily than infix notation because infix requires
parentheses but postfix does not.
\end{pastexam}

\begin{pastsoln}
答案是 (A)，對。

中序式光看符號的先後順序決定不了計算順序：$a + b * c$ 要先算乘法，是靠運算子的優先順序；要先
算加法就得寫成 $(a + b) * c$，靠括號。 所以處理中序式時，程式必須知道優先順序和結合性，還要處理
括號。 後序式的計算順序完全由符號排列決定：$(a + b) * c$ 是 \texttt{a b + c *}，$a + b * c$ 是
\texttt{a b c * +}，任何運算式都不需要括號，也不必查優先順序。

因此後序式只要從左到右掃一遍、用一個堆疊：看到運算元就推入，看到運算子就彈出兩個運算元，算完
把結果推回（資結書第~6.3.2~節）。 中序式要先轉成後序（資結書第~6.3.3~節，又要一個堆疊處理
優先順序和括號）或用兩個堆疊才能求值。 題目說的理由正是「後序式不需要括號（也就不需要優先順序
規則）」，敘述正確。

嚴格說，中序式並不是每一個都要括號（$a + b * c$ 就沒有），但有些計算順序非用括號不能表達，
後序式則永遠不需要；題目的意思是這個，不必因此判錯。
\end{pastsoln}
